% !TeX root = ../main.tex
\section{Уравнение Гельмгольца}
\label{sec:helmholtz}

Рассмотрим монохроматическое поле
с временной зависимостью
$e^{-i\omega t}$:

\begin{equation}
\mathbf{E}(\mathbf{r},t)
=
\operatorname{Re}
\left[
\widetilde{\mathbf{E}}(\mathbf{r})
e^{-i\omega t}
\right].
\end{equation}

Подстановка в волновое уравнение
для однородной среды даёт:

\begin{equation}
\nabla^2\widetilde{\mathbf{E}}
+
k^2\widetilde{\mathbf{E}}
=0,
\end{equation}

где

\begin{equation}
k=\omega\sqrt{\mu\varepsilon}.
\end{equation}

Для допустимой скалярной постановки
получаем:

\begin{equation}
\boxed{
\nabla^2 u+k^2u=0
}.
\end{equation}

Скалярное уравнение не заменяет
в общем случае векторную систему
Максвелла и требует дополнительных
предположений о среде и поляризации.
